\begin{table}[H]
\centering
\caption{Direct status-file reasons for holdout outcomes}
\label{tab:status-reasons}
\small
\begin{tabular}{lrrr}
\toprule
Recorded reason & 2019--2022 & 2023 & 2024 \\
\midrule
Repeated losses (01) & 22 & 49 & 24 \\
Net-asset deterioration (02) & 9 & 12 & 18 \\
Adverse audit opinion (03) & 111 & 23 & 37 \\
Other specified condition (04) & 2 & 1 & 1 \\
Operational impairment (06) & 87 & 21 & 25 \\
Stalled major matter (07) & 9 & 0 & 0 \\
Frozen bank account (08) & 4 & 0 & 1 \\
Information-disclosure violation (11) & 6 & 26 & 31 \\
Investment risk (12) & 9 & 2 & 1 \\
Other abnormal condition (14) & 0 & 0 & 1 \\
\midrule
Any outcome & 246 & 109 & 119 \\
\bottomrule
\end{tabular}
\begin{minipage}{0.94\textwidth}\footnotesize
\textit{Notes:} Fiscal year indexes the financial record; the designation occurs after its report landmark. One event may carry multiple reason codes, so columns need not sum to the number of outcomes. Code 11 records an information-disclosure violation. Its count rises from 6 of 246 outcomes in 2019--2022 to 57 of 228 in 2023--2024. This is a descriptive decomposition, not a causal estimate of the rule revision.
\end{minipage}
\end{table}
